% !PRINTED 597
\weberpage{597}

The totality of the integral and fractional numbers of the field $\Omega$, to
which we assign an arbitrary excludent $S$, forms the number group $O$, and the
totality of the ideals $\overline{O}$ in the same sense an \gesperrt{ideal group}.

Likewise the totality of the \gesperrt{numerical units} forms a number group $E$, and the
totality of the \gesperrt{functional units} an ideal group $\overline{E}$.

The multiplication of two groups $A$ and $B$ into the product $AB$ is effected
by multiplying every element of $A$ by every element of $B$. If $A$ is a number
group, then $\overline{E}A$ is an ideal group. But this contains only principal
ideals and can therefore be called a \gesperrt{principal ideal group}.

If the whole group $O$ splits into a \gesperrt{finite number of cosets} with
respect to $A$:
\weq{1}{O = A\alpha_1 + A\alpha_2 + \cdots + A\alpha_j,}
then the number $j$ is called the index of $A$ and is denoted by
\weq{2}{j = (O, A)}
(as in § 100).

We state the first theorem:

1. \gesperrt{If $A$ is a number group with finite index $j$, and $\eta$ any
number in $O$, then $\eta^j$ is contained in $A$.}

For of the unlimited sequence of powers $1, \eta, \eta^2, \eta^3, \dots$, two
different terms must be contained in the same coset $A\alpha_i$, and
their quotient, which is of course also a power of $\eta$, is contained in $A$.
The lowest exponent $r$ for which $\eta^r$ is contained in $A$ proves, in the
well-known way, to be a divisor of $j$.

2. \gesperrt{If $A$ is a number group with finite index, then the ideal group
$\overline{O}$ can likewise be decomposed with respect to $\overline{E}A$ into a
finite number of cosets:}
\weq{3}{\begin{aligned}
\overline{O} &= \overline{A}_1 + \overline{A}_2 + \cdots + \overline{A}_h,\\
&= A\mathfrak{a}_1 + A\mathfrak{a}_2 + \cdots + A\mathfrak{a}_h,
\end{aligned}}
and indeed we can therein choose $\mathfrak{a}_1, \mathfrak{a}_2, \dots,
\mathfrak{a}_h$ as integral ideals in $\overline{O}$, since by 1. we can
convert every fractional ideal in $\overline{O}$ into an integral ideal by
multiplication with a number in $A$. We need only choose for $\eta$ a number
that is divisible by the denominator of the ideal, and multiply by a
sufficiently high power of $\eta$.
